设总体比例为
\begin{align*}
p_{1}&=\text{在谢尔伍德银行存款的储户的比例}\\
p_{2}&=\text{在银行 B 存款的储户的比例}\\
p_{3}&=\text{在银行 C 存款的储户的比例}\\
p_{4}&=\text{在银行 D 存款的储户的比例}\\
1-(p_{1}&+p_{2}+p_{3}+p_{4})=\text{在其他银行存款的储户的比例}
\end{align*}

(a) 求 $p_{1},p_{2},\cdots,p_{4}$ 的 $95\%$ 联合置信区间

(b) 构造一个可将谢尔伍德银行与其主要竞争对手 B 进行比较的 $95\%$ 联合置信区间, 并对这一区间作出解释.

\medskip
\noindent\mvbold{5.17}\quad 为评估某城市高中学生中吸毒问题的普遍性, 对来自该城市 5 所高中的 200 名学生所组成的一个随机样本进行了调查, 调查的问题之一及相应的回答如下:

通常您一周吸几支大麻烟?

\begin{mvtable}
\centering
\begin{tabular}{|c|c|c|c|}
\hline
\multirow{2}{*}{} & \multicolumn{3}{c|}{分类}\\
\cline{2-4}
 & 不吸 & 中等 ($1\sim3$ 支) & 严重 ($4$ 支以上)\\
\hline
应答人数 & $117$ & $62$ & $21$\\
\hline
\end{tabular}
\end{mvtable}

试对这三种人的比例 $p_{1},p_{2}$ 与 $p_{3}=1-(p_{1}+p_{2})$ 构造 $95\%$ 联合置信区间.

\medskip
\noindent\mvbold{下面的练习需要使用计算机.}

\medskip
\noindent\mvbold{5.18}\quad 利用\hyperref[tab:5.2]{表 5.2} 中的大学考试数据 (见例 5.5)

\noindent (a) 在显著性水平 $\alpha=0.05$ 下检验假设
\[
H_{0}:\boldsymbol{\mu}'=[500,50,30],\quad H_{1}:\boldsymbol{\mu}'\ne[500,50,30]
\]
假定 $[500,50,30]'$ 代表近 10 年来数千名大学生考试的平均成绩. 是否有理由相信, 成绩被列入\hyperref[tab:5.2]{表 5.2} 的那一组学生与以往的学生相比其得分情况有明显不同? 请作出解释.

\noindent (b) 确定 $\boldsymbol{\mu}$ 的 $95\%$ 置信椭球的主轴的长度与方向.

\noindent (c) 从社会科学与历史、语言表达以及自然科学各科成绩的边缘分布画出 Q-Q 图. 从各对不同变量的观测值中画出可能的三个散布图. 讨论这些数据是否服从正态分布.

\medskip
\noindent\mvbold{5.19}\quad \hyperref[tab:5.11]{表 5.11} 中给出由 $n=30$ 根一等硬木组成的样本的硬度 $x_{1}$ 与弯曲强度 $x_{2}$ 的测量值, 单位为磅/英寸$^{2}$. 利用表中数据:

\begin{mvtable}
\phantomsection\label{tab:5.11}\mvtabcap{表 5.11\quad 木材数据}
\centering
\begin{tabular}{|c|c||c|c|}
\hline
$x_{1}$(硬度) & $x_{2}$(弯曲强度) & $x_{1}$(硬度) & $x_{2}$(弯曲强度)\\
\hline
$1\,232$ & $4\,175$ & $1\,712$ & $7\,749$\\
$1\,115$ & $6\,652$ & $1\,932$ & $6\,818$\\
$2\,205$ & $7\,612$ & $1\,820$ & $9\,307$\\
$1\,897$ & $10\,914$ & $1\,900$ & $6\,457$\\
$1\,932$ & $10\,850$ & $2\,426$ & $10\,102$\\
$1\,612$ & $7\,627$ & $1\,558$ & $7\,414$\\
$1\,598$ & $6\,954$ & $1\,470$ & $7\,556$\\
$1\,804$ & $8\,365$ & $1\,858$ & $7\,833$\\
\hline
\end{tabular}
\end{mvtable}
